\justify
La variabilidad de precios en los materiales de construcción no es un fenómeno exclusivo de Paraguay. En buena parte de América Latina, la dependencia de materiales importados y la exposición a shocks externos (cambiarios, energéticos y sanitarios) generan fluctuaciones que afectan a toda la cadena del sector. Sin embargo, cuando se revisa la literatura regional, lo que predomina son diagnósticos de lo ocurrido, no herramientas para anticiparlo. La mayoría de los estudios adoptan enfoques descriptivos o correlacionales; los que incorporan modelos de series temporales o aprendizaje profundo son todavía escasos.
\justify
En Colombia, Castañeda y Aguirre (2022) analizaron la variabilidad de precios de materiales e insumos para la construcción en el departamento del Valle del Cauca \cite{castaneda2022}. A partir de los índices de costos de construcción de vivienda proporcionados por la Cámara Colombiana de la Construcción (CAMACOL) para el período enero de 2019 a agosto de 2021, los autores llevaron a cabo una investigación cuantitativa de carácter exploratorio-descriptivo. Sus resultados mostraron que la pandemia de COVID-19, las protestas sociales de 2021 y el encarecimiento de insumos importados fueron los principales factores detrás del alza de precios, con incrementos de hasta un 33,5\,\% en el acero durante 2021. El estudio también señaló que estas variaciones provocaron retrasos, modificaciones presupuestarias y sobrecostos en la ejecución de obras, y destacó el papel de la interventoría técnica como figura relevante para mitigar los impactos económicos de la inflación de materiales.
\justify
En Paraguay, Almada de Bernal et al.\ (2024) estudiaron la variación de precios de materiales e insumos para la construcción en Ciudad del Este \cite{almada2022}. Los autores relevaron precios en tres comercios locales entre octubre de 2021 y octubre de 2022, comparando las variaciones porcentuales de 100 productos de construcción civil, entre ellos cementos, ladrillos, varillas de hierro, pisos y caños de PVC. El análisis evidenció subas y bajas dispares según el comercio y el producto, con variaciones que llegaron al 245\,\% en algunos tipos de varillas de hierro. Los autores atribuyeron estos movimientos al efecto pospandemia y a la adaptación de las empresas fabricadoras y distribuidoras a los precios globales, según datos de la Cámara Paraguaya de Construcción (CAPACO).
\justify
Frente a ese panorama, el presente trabajo toma una dirección distinta. En lugar de describir variaciones ya ocurridas, se propone construir modelos capaces de proyectarlas. Para ello se combinan métodos econométricos (SARIMAX) con arquitecturas de redes neuronales recurrentes (LSTM y GRU), e incorpora como variables externas tanto el nivel del río Paraguay, un factor hidrológico con incidencia logística directa, como el confinamiento por COVID-19. El objetivo no es solo medir la precisión de los modelos, sino generar una herramienta concreta que ayude a tomar mejores decisiones en el sector de la construcción paraguayo.
